% GENERATED FILE. Edit canonical lessons, then run npm run generate:books.
% canonical-source-hash: fnv1a64:aa180b5c9787c6d5
% canonical-lessons: KA-C428-ondomme, KA-C428-illave, KA-C428-hegiddaru, KA-C428-anvaya, KA-C428-ottare

\chapter{If and Otherwise}
\label{ch:ka-if-and-otherwise}
\hlchaptermodality{428}

\begin{chapteropening}
\textbf{By the end of this chapter:} \emph{I can say if ever, or else, in any case, applicability and overall.}

Complete the last lesson of chapter 428: I can say if ever, or else, in any case, applicability and overall.
\end{chapteropening}

\section[\texorpdfstring{\kn{ಒಂದೊಮ್ಮೆ} (ondomme)}{ಒಂದೊಮ್ಮೆ (ondomme)}]{\kn{ಒಂದೊಮ್ಮೆ} (ondomme) — if ever; should it happen}

\label{lesson:KA-C428-ondomme}



\begin{quote}
Before the new one: say the Kannada for to convert, then the Kannada for adequate.
\end{quote}

\subsection*{You'll want to know: \kn{ಒಂದೊಮ್ಮೆ}}
\textbf{\kn{ಒಂದೊಮ್ಮೆ}} — \emph{ondomme} — \textquotedblleft{}if ever; should it happen\textquotedblright{}.

\begin{quote}
\kn{ಒಂದೊಮ್ಮೆ} \kn{ಬಸ್} \kn{ತಪ್ಪಿದರೆ} \kn{ನಾನು} \kn{ಆಟೋದಲ್ಲಿ} \kn{ಬರುತ್ತೇನೆ}.
\end{quote}

\begin{quote}
\emph{ondomme bas tappidare nānu āṭōdalli baruttēne.}
\end{quote}

\begin{quote}
Should I miss the bus, I will come by auto.
\end{quote}

It goes with a verb in -\kn{ದರೆ}, as \kn{ಒಂದು} \kn{ವೇಳೆ} does.

\begin{cousinweb}
\textbf{\kn{ಒಂದೊಮ್ಮೆ}} (\emph{ondomme}) is native Dravidian: \emph{ondu}, \textquotedblleft{}one\textquotedblright{}, and \emph{omme}, \textquotedblleft{}once\textquotedblright{} (DEDR 990): at some one time.
\end{cousinweb}

\begin{grammarlens}[title={What you've built}]
One more word for saying what happens if.
\end{grammarlens}

\subsection*{Your turn}
\begin{itemize}
\raggedright
  \item \emph{Say it:} \emph{ondomme}
  \item \emph{Say it:} \emph{ondomme}, once more
  \item \emph{Say it:} \emph{ondomme bas tappidare nānu āṭōdalli baruttēne.} — Should I miss the bus, I will come by auto.
  \item \emph{Recall:} say the Kannada for to convert, then the Kannada for adequate, then say \emph{ondomme} again
\end{itemize}

\begin{tcolorbox}[breakable,colback=teal!4,colframe=teal!35!black,title={Before you move on}]
What does \kn{ಒಂದೊಮ್ಮೆ} mean? (\textquotedblleft{}If ever; should it happen\textquotedblright{}.) Say it once more.
\end{tcolorbox}
\section[\texorpdfstring{\kn{ಇಲ್ಲವೇ} (illavē)}{ಇಲ್ಲವೇ (illavē)}]{\kn{ಇಲ್ಲವೇ} (illavē) — or else; or}

\label{lesson:KA-C428-illave}



\begin{quote}
Before the new one: say the Kannada for adequate, then the Kannada for if ever.
\end{quote}

\subsection*{You'll want to know: \kn{ಇಲ್ಲವೇ}}
\textbf{\kn{ಇಲ್ಲವೇ}} — \emph{illavē} — \textquotedblleft{}or else; or\textquotedblright{}.

\begin{quote}
\kn{ಇಂದು} \kn{ಇಲ್ಲವೇ} \kn{ನಾಳೆ} \kn{ಮಳೆ} \kn{ಬರುತ್ತದೆ}.
\end{quote}

\begin{quote}
\emph{indu illavē nāḷe maḷe baruttade.}
\end{quote}

\begin{quote}
It will rain today or tomorrow.
\end{quote}

Like \kn{ಅಥವಾ}, but stronger: one or the other.

\begin{cousinweb}
\textbf{\kn{ಇಲ್ಲವೇ}} (\emph{illavē}) is native Dravidian: \emph{illa}, \textquotedblleft{}it is not\textquotedblright{}, with -\emph{ē}; DEDR 2559 lists the Kannada family of \emph{illa}.
\end{cousinweb}

\begin{grammarlens}[title={What you've built}]
One more word for saying what happens if.
\end{grammarlens}

\subsection*{Your turn}
\begin{itemize}
\raggedright
  \item \emph{Say it:} \emph{illavē}
  \item \emph{Say it:} \emph{illavē}, once more
  \item \emph{Say it:} \emph{indu illavē nāḷe maḷe baruttade.} — It will rain today or tomorrow.
  \item \emph{Recall:} say the Kannada for adequate, then the Kannada for if ever, then say \emph{illavē} again
\end{itemize}

\begin{tcolorbox}[breakable,colback=teal!4,colframe=teal!35!black,title={Before you move on}]
What does \kn{ಇಲ್ಲವೇ} mean? (\textquotedblleft{}Or else; or\textquotedblright{}.) Say it once more.
\end{tcolorbox}
\section[\texorpdfstring{\kn{ಹೇಗಿದ್ದರೂ} (hēgiddarū)}{ಹೇಗಿದ್ದರೂ (hēgiddarū)}]{\kn{ಹೇಗಿದ್ದರೂ} (hēgiddarū) — in any case; anyway}

\label{lesson:KA-C428-hegiddaru}



\begin{quote}
Before the new one: say the Kannada for if ever, then the Kannada for or else.
\end{quote}

\subsection*{You'll want to know: \kn{ಹೇಗಿದ್ದರೂ}}
\textbf{\kn{ಹೇಗಿದ್ದರೂ}} — \emph{hēgiddarū} — \textquotedblleft{}in any case; anyway\textquotedblright{}.

\begin{quote}
\kn{ಹೇಗಿದ್ದರೂ} \kn{ನಾನು} \kn{ಪೇಟೆಗೆ} \kn{ಹೋಗುತ್ತೇನೆ}, \kn{ನಿನ್ನ} \kn{ಪುಸ್ತಕ} \kn{ತರುತ್ತೇನೆ}.
\end{quote}

\begin{quote}
\emph{hēgiddarū nānu pēṭege hōguttēne, ninna pustaka taruttēne.}
\end{quote}

\begin{quote}
I am going to town anyway, so I will bring your book.
\end{quote}

\kn{ಹೇಗೆ} (how) + \kn{ಇದ್ದರೂ} (even if it is).

\begin{cousinweb}
\textbf{\kn{ಹೇಗಿದ್ದರೂ}} (\emph{hēgiddarū}) is native Dravidian: \emph{hēge}, \textquotedblleft{}how\textquotedblright{} (DEDR 5151), and \emph{iddarū}, \textquotedblleft{}even if it is\textquotedblright{}: however it is.
\end{cousinweb}

\begin{grammarlens}[title={What you've built}]
One more word for saying what happens if.
\end{grammarlens}

\subsection*{Your turn}
\begin{itemize}
\raggedright
  \item \emph{Say it:} \emph{hēgiddarū}
  \item \emph{Say it:} \emph{hēgiddarū}, once more
  \item \emph{Say it:} \emph{hēgiddarū nānu pēṭege hōguttēne, ninna pustaka taruttēne.} — I am going to town anyway, so I will bring your book.
  \item \emph{Recall:} say the Kannada for if ever, then the Kannada for or else, then say \emph{hēgiddarū} again
\end{itemize}

\begin{tcolorbox}[breakable,colback=teal!4,colframe=teal!35!black,title={Before you move on}]
What does \kn{ಹೇಗಿದ್ದರೂ} mean? (\textquotedblleft{}In any case; anyway\textquotedblright{}.) Say it once more.
\end{tcolorbox}
\section[\texorpdfstring{\kn{ಅನ್ವಯ} (anvaya)}{ಅನ್ವಯ (anvaya)}]{\kn{ಅನ್ವಯ} (anvaya) — applicability}

\label{lesson:KA-C428-anvaya}



\begin{quote}
Before the new one: say the Kannada for or else, then the Kannada for in any case.
\end{quote}

\subsection*{You'll want to know: \kn{ಅನ್ವಯ}}
\textbf{\kn{ಅನ್ವಯ}} — \emph{anvaya} — \textquotedblleft{}applicability\textquotedblright{}.

\begin{quote}
\kn{ಈ} \kn{ನಿಯಮ} \kn{ಎಲ್ಲರಿಗೂ} \kn{ಅನ್ವಯ} \kn{ಆಗುತ್ತದೆ}.
\end{quote}

\begin{quote}
\emph{ī niyama ellarigū anvaya āguttade.}
\end{quote}

\begin{quote}
This rule applies to everyone.
\end{quote}

\kn{ಅನ್ವಯ} \kn{ಆಗು} / \kn{ಅನ್ವಯಿಸು}, to apply to.

\begin{cousinweb}
\textbf{\kn{ಅನ್ವಯ}} (\emph{anvaya}) is borrowed from Sanskrit \emph{anvaya}, \textquotedblleft{}connection, being linked to or concerned with\textquotedblright{} (Monier-Williams).
\end{cousinweb}

\begin{grammarlens}[title={What you've built}]
One more word for saying what happens if.
\end{grammarlens}

\subsection*{Your turn}
\begin{itemize}
\raggedright
  \item \emph{Say it:} \emph{anvaya}
  \item \emph{Say it:} \emph{anvaya}, once more
  \item \emph{Say it:} \emph{ī niyama ellarigū anvaya āguttade.} — This rule applies to everyone.
  \item \emph{Recall:} say the Kannada for or else, then the Kannada for in any case, then say \emph{anvaya} again
\end{itemize}

\begin{tcolorbox}[breakable,colback=teal!4,colframe=teal!35!black,title={Before you move on}]
What does \kn{ಅನ್ವಯ} mean? (\textquotedblleft{}Applicability\textquotedblright{}.) Say it once more.
\end{tcolorbox}
\section[\texorpdfstring{\kn{ಒಟ್ಟಾರೆ} (oṭṭāre)}{ಒಟ್ಟಾರೆ (oṭṭāre)}]{\kn{ಒಟ್ಟಾರೆ} (oṭṭāre) — overall; on the whole}

\label{lesson:KA-C428-ottare}



\begin{quote}
Before the new one: say the Kannada for in any case, then the Kannada for applicability.
\end{quote}

\subsection*{You'll want to know: \kn{ಒಟ್ಟಾರೆ}}
\textbf{\kn{ಒಟ್ಟಾರೆ}} — \emph{oṭṭāre} — \textquotedblleft{}overall; on the whole\textquotedblright{}.

\begin{quote}
\kn{ಒಟ್ಟಾರೆ} \kn{ಪ್ರವಾಸ} \kn{ಚೆನ್ನಾಗಿತ್ತು}.
\end{quote}

\begin{quote}
\emph{oṭṭāre pravāsa cennāgittu.}
\end{quote}

\begin{quote}
Overall the trip was good.
\end{quote}

Built on \kn{ಒಟ್ಟು}, the total.

\begin{cousinweb}
\textbf{\kn{ಒಟ್ಟಾರೆ}} (\emph{oṭṭāre}) is native Dravidian, built on \emph{oṭṭu}, \textquotedblleft{}to come together; a mass, a union\textquotedblright{} (DEDR 958).
\end{cousinweb}

\begin{grammarlens}[title={What you've built}]
One more word for saying what happens if.
\end{grammarlens}

\subsection*{Your turn}
\begin{itemize}
\raggedright
  \item \emph{Say it:} \emph{oṭṭāre}
  \item \emph{Say it:} \emph{oṭṭāre}, once more
  \item \emph{Say it:} \emph{oṭṭāre pravāsa cennāgittu.} — Overall the trip was good.
  \item \emph{Recall:} say the Kannada for in any case, then the Kannada for applicability, then say \emph{oṭṭāre} again
\end{itemize}

\begin{tcolorbox}[breakable,colback=teal!4,colframe=teal!35!black,title={Before you move on}]
What does \kn{ಒಟ್ಟಾರೆ} mean? (\textquotedblleft{}Overall; on the whole\textquotedblright{}.) Say it once more.
\end{tcolorbox}
